\documentclass[a4paper,10pt]{article}
\usepackage[francais]{babel}
\usepackage{url}

\selectlanguage{francais}
\input{letter-ensiie.sty}

\selectlanguage{francais}
\enCouleur

\begin{document}


\name{St\'ephane \textsc{Genaud}}
\fonction{Professeur des Universités, ENSIIE}
\telephoneillkirch{+33 (0)3688 54542}
\email{genaud@ensiie.fr}

\begin{letter}[Recommendation de Slimane Agouram]{à}%
{M. Robin Ryder}


Cher Collègue,


Slimane Agouram  a obtenu le diplôme  d'ingénieur de l'ENSIIE en  novembre 2015,
après une scolarité de 3 ans. La formation d'ingénieur de l'ENSIIE est dispensée
à  la fois  sur  le site  historique  principal, à  Evry, et  à  son antenne  de
Strasbourg.   Elle   est  composée   d'un  tronc  commun   puis  d'enseignements
spécialisés.   Le  contenu  pédagogique  du  tronc  commun  de  trois  semestres
(contenus pédagogiques et  examens communs aux deux sites) est  constitué à 40\%
d'informatique,  40\%   de  mathématiques,   et  20\%  d'enseignement   de  type
socio-économique. La fin du cursus est composée, en plus des enseignements
généraux, d'une spécialité suivie dans le cadre d'un Master de l'université de
Strasbourg en bi-cursus. Dans ce contexte, Slimane a obtenu un Master en
informatique, parcours \emph{Ingénierie Logicielle et Connaissances}, abordant
en particulier la fouille de données.

En tant  que directeur de  l'antenne de Strasbourg,  j'ai suivi Slimane  tout au
long de ses  trois années de scolarité, parfois en  ayant Slimane comme étudiant
dans certains cours.  Je peux  témoigner de la personnalité extrêmement positive
de Slimane. Il a été tout long de sa formation, un étudiant curieux, attentif et
travailleur. Ceci lui  a permis de valider  sa formation de manière  tout à fait
satisfaisante. Je suis également convaincu que les qualités humaines de Slimane
lui permettront de poursuivre une belle carrière professionnelle.

Depuis son  entrée dans la vie  professionnelle, Slimane s'est orienté  vers les
problématiques Big  Data. Il travaille  dans ce  domaine en tant  que consultant
décisionnel  BI/Big Data  depuis  deux ans.  Dans un  entretien  récent, il  m'a
expliqué son grand  intérêt pour la science  des données, qui se  traduit par la
lecture d'ouvrages sur le sujet.

Pour  ces raisons,  je suis  convaincu de  sa motivation  pour suivre  le Master
executive   \emph{Statistique  et   Big  Data}   enseignée  à   l'université  de
Paris-Dauphine, et je ne peux que vous recommander de lui donner cette chance.

Bien cordialement,

\end{letter}
\begin{flushright}
%\includegraphics[width=.2\textwidth]{signgenaud.jpg}
\end{flushright}
\end{document}
